\documentclass[11pt,a4paper]{article}
\usepackage[spanish,es-noquoting]{babel}
\usepackage[top=2cm, bottom=2cm, left=2cm, right=2cm]{geometry}
\usepackage{fontspec}
\usepackage{amsmath,amssymb}
\usepackage[version=4]{mhchem}
\usepackage{enumitem}
\usepackage{microtype}

\setmainfont{Calibri}
\linespread{1.15}

\begin{document}
\pagestyle{empty}

\begin{center}
    {\LARGE \textbf{Bolsa de Ejercicios — 4.º ESO / 1.º de Bachillerato}}\\[6pt]
    {\normalsize Nombre: \underline{\hspace{8cm}} \quad Fecha: \underline{\hspace{3cm}}}
\end{center}

\vspace{0.5cm}
\noindent\textbf{1.} Un futbolista golpea un balón situado en reposo sobre el césped horizontal comunicándole una velocidad inicial $v_0 = 20\text{ m/s}$ con un ángulo de elevación de $\alpha = 30^\circ$ sobre el suelo horizontal. Despreciando el rozamiento con el aire y tomando $g = 9{,}8\text{ m/s}^2$:
\begin{enumerate}[label=\textbf{\alph*)}, leftmargin=2em, itemsep=6pt]
    \item Escribe las ecuaciones del movimiento en cada eje ($X$ e $Y$).
    \item Calcula el tiempo que tarda en alcanzar el punto más alto de su trayectoria y determina la altura máxima alcanzada sobre el suelo.
    \item Determina el tiempo total de permanencia en el aire hasta el instante en que contacta de nuevo contra el terreno de juego.
    \item Calcula la distancia horizontal medida sobre el suelo entre el punto inicial de lanzamiento y el punto de impacto.
    \item Halla las componentes cartesianas del vector velocidad y el módulo de la velocidad en el instante inmediatamente anterior a la colisión con el suelo.
\end{enumerate}

\vspace{0.4cm}
\noindent\textbf{2.} Desde el borde superior de un acantilado vertical situado a una altura $h = 45\text{ m}$ sobre el nivel del mar, una dotación de emergencia lanza una cápsula de salvamento marítimo con una velocidad inicial $v_0 = 30\text{ m/s}$ orientada a un ángulo de $\alpha = 30^\circ$ por encima de la línea horizontal. Se considera despreciable la resistencia aerodinámica y se toma $g = 9{,}8\text{ m/s}^2$.
\begin{enumerate}[label=\textbf{\alph*)}, leftmargin=2em, itemsep=6pt]
    \item Escribe las ecuaciones del movimiento en cada eje ($X$ e $Y$).
    \item Halla el tiempo que transcurre hasta alcanzar la cota más elevada del vuelo y determina dicha altura máxima absoluta respecto a la superficie marina.
    \item Determina el tiempo total que tarda la cápsula desde el instante de su lanzamiento hasta que impacta contra el agua del mar ($y = 0$).
    \item Calcula la distancia horizontal comprendida entre la base vertical del acantilado y el punto exacto de entrada en el agua marina.
    \item Calcula las componentes del vector velocidad de entrada en el agua y halla el módulo de la velocidad final de impacto.
\end{enumerate}

\newpage

\noindent\textbf{3.} Sobre una rampa rectilínea que forma un ángulo $\alpha = 30^\circ$ con el suelo horizontal descansa un cajón metálico de masa $m = 8\text{ kg}$. El coeficiente de rozamiento dinámico entre el cajón y la superficie de la rampa es $\mu = 0{,}20$. Para desplazarlo pendiente arriba, se aplica una fuerza constante $F = 85\text{ N}$ paralela en todo momento a la superficie de la rampa y orientada en sentido ascendente ($g = 9{,}8\text{ m/s}^2$).
\begin{enumerate}[label=\textbf{\alph*)}, leftmargin=2em, itemsep=6pt]
    \item Dibuja el diagrama de cuerpo libre detallado, identificando y rotulando todas las fuerzas actuantes ($\vec{P}$, $\vec{N}$, $\vec{F}_r$ y $\vec{F}$) y descomponiendo vectorialmente el peso en las direcciones paralela y perpendicular al plano inclinado.
    \item Calcula el valor numérico exacto de la fuerza normal ejercida por el plano ($N$) y de la fuerza de rozamiento por deslizamiento ($F_r$).
    \item Aplica el principio fundamental de la dinámica (segunda ley de Newton) para deducir la aceleración $a$ con la que asciende el cajón.
    \item Sabiendo que el cajón arranca desde el reposo ($v_0 = 0$), halla el tiempo requerido para avanzar una distancia de $10\text{ m}$ a lo largo de la rampa y calcula la velocidad final que alcanza en ese instante.
\end{enumerate}

\vspace{0.4cm}
\noindent\textbf{4.} En un ensayo de laboratorio se dispone de una probeta con $250\text{ mL}$ de una disolución acuosa de ácido sulfúrico (\ce{H2SO4}) de concentración $0{,}40\text{ M}$. Se desea neutralizar completamente dicho volumen mediante la adición controlada de hidróxido de sodio (\ce{NaOH}).
\begin{enumerate}[label=\textbf{\alph*)}, leftmargin=2em, itemsep=6pt]
    \item Escribe y ajusta la ecuación química de la reacción de neutralización.
    \item Calcula la cantidad de sustancia (en moles) de ácido sulfúrico presentes en la muestra inicial y los moles de hidróxido de sodio estequiométricamente precisos para neutralizarlo.
    \item Determina la masa en gramos de reactivo puro \ce{NaOH} correspondiente a dicha cantidad.
    \item Si en el laboratorio se dispone de una disolución comercial de sosa acuosa $1{,}25\text{ M}$, ¿qué volumen en mililitros de esta disolución se debe pipetear para completar la neutralización?
    \item Calcula la masa en gramos de sulfato de sodio (\ce{Na2SO4}) que se obtendrá en la disolución al concluir la reacción.
\end{enumerate}
\noindent\textbf{Datos atómicos:} Masas atómicas relativas ($\text{u}$): $\mathrm{H} = 1{,}0;\; \mathrm{O} = 16{,}0;\; \mathrm{Na} = 23{,}0;\; \mathrm{S} = 32{,}0$.

\vspace{0.4cm}
\noindent\textbf{5.} En un reactor industrial se introducen $56{,}0\text{ g}$ de gas nitrógeno puro (\ce{N2}) y $18{,}0\text{ g}$ de gas hidrógeno puro (\ce{H2}) a presión elevada para sintetizar amoniaco gaseoso según la reacción sin ajustar: $\ce{N2 (g) + H2 (g) -> NH3 (g)}$.
\begin{enumerate}[label=\textbf{\alph*)}, leftmargin=2em, itemsep=6pt]
    \item Ajusta la ecuación química y explica el significado molecular y molar de los coeficientes estequiométricos.
    \item Determina de forma cuantitativa y razonada cuál de los dos reactivos actúa como reactivo limitante y cuál está en exceso.
    \item Calcula la masa en gramos del reactivo en exceso que quedará sin consumir al agotarse el reactivo limitante.
    \item Calcula la masa teórica máxima de gas amoniaco (\ce{NH3}) susceptible de formarse suponiendo un rendimiento del $100\%$.
    \item Si las condiciones reales de operación fijan un rendimiento de la reacción del $85{,}0\%$, calcula la masa real de amoniaco producida y el volumen real que ocupará medido a temperatura ambiente de $25^\circ\text{C}$ ($298\text{ K}$) y $1\text{ atm}$ de presión ($R = 0{,}082\text{ atm}\cdot\text{L}/(\text{mol}\cdot\text{K})$).
\end{enumerate}
\noindent\textbf{Constantes y datos:} Masas atómicas relativas ($\text{u}$): $\mathrm{H} = 1{,}0;\; \mathrm{N} = 14{,}0$. Constante de los gases ideales: $R = 0{,}082\text{ atm}\cdot\text{L}/(\text{mol}\cdot\text{K})$ ($V_m \approx 24{,}4\text{ L/mol}$ a $25^\circ\text{C}$ y $1\text{ atm}$).

\newpage

\noindent\textbf{6.} Se quiere enlosar un jardín con forma de triángulo rectángulo de catetos $42\text{ m}$ y $56\text{ m}$. Dentro del jardín se va a diferenciar un espacio rectangular techado de forma que dos de sus lados sean paralelos a los catetos del triángulo, un vértice coincida con el vértice del ángulo recto del triángulo y el vértice opuesto esté sobre su hipotenusa. Alicatar la parte cubierta cuesta $30\text{ €/m}^2$ y la parte no techada, $50\text{ €/m}^2$ pues las baldosas llevan un tratamiento especial resistente al agua. Calcule las dimensiones de la parte techada que hacen que el coste de instalar el suelo en el jardín sea mínimo.

\vspace{0.4cm}
\noindent\textbf{7.} Se considera la función real de variable real definida por la expresión analítica:
\[
f(x) = \frac{|2x - 4|}{x + 1}
\]
\begin{enumerate}[label=\textbf{\alph*)}, leftmargin=2em]
    \item Pinta la función.
\end{enumerate}

\end{document}
